\section*{Hoofdstuk \ref{ch:g4:clouds-rain-snow} --- Wolken, regen en sneeuw}

\begin{solution}{exo:g4:clouds-rain-snow:1}
Een \omterm{def:g4:clouds-rain-snow:cloud}{wolk} bestaat uit een enorm aantal piepkleine waterdruppeltjes, of
piepkleine stukjes ijs, die in de lucht zweven. \omterm{def:g4:clouds-rain-snow:cloud}{Mist} is een \omterm{def:g4:clouds-rain-snow:cloud}{wolk} die de
grond raakt.
\end{solution}

\begin{solution}{exo:g4:clouds-rain-snow:2}
(a) \omterm{def:g4:clouds-rain-snow:cloud-families}{Stratus}. (b) \omterm{def:g4:clouds-rain-snow:cloud-families}{Cirrus}. (c) \omterm{def:g4:clouds-rain-snow:cloud-families}{Cumulus}. (d) \omterm{def:g4:clouds-rain-snow:cloud-families}{Cumulonimbus}.
\end{solution}

\begin{solution}{exo:g4:clouds-rain-snow:3}
(a) \omterm{def:g4:clouds-rain-snow:precipitation}{Hagel}. (b) \omterm{def:g4:clouds-rain-snow:precipitation}{Regen}. (c) \omterm{def:g4:clouds-rain-snow:precipitation}{Sneeuw}.
\end{solution}

\begin{solution}{exo:g4:clouds-rain-snow:4}
(a) Een \omterm{def:g4:clouds-rain-snow:rain-gauge}{regenmeter}. (b) Een thermometer. (c) Een \omterm{def:g4:clouds-rain-snow:wind-vane}{windvaan}.
\end{solution}

\begin{solution}{exo:g4:clouds-rain-snow:5}
Er is 17 millimeter \omterm{def:g4:clouds-rain-snow:precipitation}{regen} gevallen. $20 - 17 = 3$: met 3 millimeter erbij
zou het water tot het streepje 20 komen.
\end{solution}

\begin{solution}{exo:g4:clouds-rain-snow:6}
De natste dag was woensdag (12 millimeter). Op maandag en op vrijdag
regende het niet. Donderdag had 7 millimeter.
\end{solution}

\begin{solution}{exo:g4:clouds-rain-snow:7}
De sneeuwvlokken ontstaan hoog in de koude \omterm{def:g4:clouds-rain-snow:cloud}{wolk}. Boven in het bergdorp
vriest het helemaal tot beneden, dus komen ze aan als \omterm{def:g4:clouds-rain-snow:precipitation}{sneeuw}. Lager, in het
\omterm{def:g1:landscapes-around-us:landform}{dal}, is de lucht warmer dan \SI{0}{\celsius}: op weg naar beneden smelten de
vlokken in die lucht, en ze bereiken de grond als \omterm{def:g4:clouds-rain-snow:precipitation}{regen}.
\end{solution}

\begin{solution}{exo:g4:clouds-rain-snow:8}
Het staafdiagram heeft zeven staven met hoogten 3, 0, 1, 9, 6, 0 en 1 (geen
staaf op dinsdag en zaterdag). In de week viel er
$3 + 0 + 1 + 9 + 6 + 0 + 1 = 20$ millimeter. Donderdag bracht 9 van de
20 millimeter: $\frac{9}{20}$ van de \omterm{def:g4:clouds-rain-snow:precipitation}{regen} van de week.
\end{solution}

\begin{solution}{exo:g4:clouds-rain-snow:9}
Muren en bomen beschutten de grond in hun buurt tegen de \omterm{def:g4:clouds-rain-snow:precipitation}{regen}, vooral als
de wind de \omterm{def:g4:clouds-rain-snow:precipitation}{regen} schuin opzij blaast, en druppels kunnen van bladeren in de
\omterm{def:g4:clouds-rain-snow:rain-gauge}{regenmeter} druipen: dichtbij zou de \omterm{def:g4:clouds-rain-snow:rain-gauge}{regenmeter} te weinig of te veel
opvangen. De thermometer moet de temperatuur van de lucht meten; in de
volle zon zou hij door het zonlicht zelf worden opgewarmd en te veel
aanwijzen. De witte hut met latjes houdt de zon tegen en laat de lucht door.
\end{solution}

\begin{solution}{exo:g4:clouds-rain-snow:10}
Een \omterm{def:g4:clouds-rain-snow:wind-vane}{windvaan} wijst naar waar de wind vandaan komt: de wind komt uit het
westen. Hij gaat naar het oosten.
\end{solution}

\begin{solution}{exo:g4:clouds-rain-snow:11}
Eén millimeter \omterm{def:g4:clouds-rain-snow:precipitation}{regen} op één vierkante meter is één liter. Op één vierkante
meter geeft 8 millimeter dus 8 liter. Op 500 vierkante meter:
$500 \times 8 = 4000$ liter. Met gieters van 10 liter:
$4000 \div 10 = 400$ gieters.
\end{solution}
